\documentclass[10pt,a4paper]{article}
\usepackage[T1]{fontenc}
\usepackage[utf8]{inputenc}
\usepackage[swedish]{babel}
\usepackage{amsmath,amssymb}
\usepackage{booktabs}
\usepackage[margin=2.0cm]{geometry}
\usepackage{microtype}

\setlength{\parindent}{0pt}
\setlength{\parskip}{0.45em}
\newcommand{\modelname}{\texttt{gerlee-history}}

\begin{document}

\begin{center}
  {\Large\bfseries Teknisk beskrivning av en adaptiv historisk regressionsmodell}\par
  \vspace{0.3em}
  {\modelname{} för den svenska prognoshubben för influensa}\par
  \vspace{0.2em}
  Philip Gerlee, Chalmers tekniska högskola
\end{center}

\section*{Syfte och datakälla}

Modellen ger punktprognoser för antalet veckovis rapporterade influensafall i
Sverige (SE), Region Skåne (SE-M) och Västra Götalandsregionen (SE-O).
Prognoser lämnas för horisonterna $h=0,1,2,3$, där $h=0$ avser veckan som
slutar på prognosens referensdatum. Modellen använder enbart den tidsserie som
publiceras av prognoshubben. Inga externa epidemiologiska eller demografiska
kovariater ingår.

Grundantagandet är att influensa uppvisar återkommande säsongsformer. Antalet
fall en viss ISO-vecka under den aktuella säsongen kan därför approximeras med
en linjär kombination av antalet fall under motsvarande vecka i tidigare
säsonger. Kombinationen anpassas successivt när nya observationer blir
tillgängliga under den aktuella säsongen.

En influensasäsong definieras från ISO-vecka 30 år $a$ till och med vecka 23
år $a+1$. Endast historiska säsonger med fullständiga observationer för alla
tre platser används. Säsongen 2020/21 utesluts eftersom pandemin och dess
smittskyddsåtgärder gjorde influensacirkulationen starkt avvikande. Vid den
retrospektiva prognosen för 2025/26 används åtta säsonger: 2016/17--2019/20
samt 2021/22--2024/25.

\section*{Modellformulering}

Betrakta en plats separat. Låt $y_t$ vara observerat antal fall vecka $t$
under den aktuella säsongen och låt
\[
  \mathbf{x}_t=(x_{t1},\ldots,x_{tn})^{\mathsf T}
\]
innehålla antalet fall under motsvarande ISO-vecka i de $n$ historiska
säsongerna. Modellens väntevärde är
\begin{equation}
  \widehat y_t = \mathbf{x}_t^{\mathsf T}\boldsymbol\beta,
  \qquad \boldsymbol\beta=(\beta_1,\ldots,\beta_n)^{\mathsf T},
  \qquad \beta_j\geq 0.
  \label{eq:forecast}
\end{equation}
I början av säsongen tilldelas alla historiska säsonger samma vikt,
\begin{equation}
  \boldsymbol\beta_0=\frac{1}{n}\mathbf{1}.
  \label{eq:prior}
\end{equation}
Modellen börjar uppdatera vikterna med observationer från vecka 30. För ett
prognosdatum $r$ får mängden $\mathcal T_r$ innehålla alla observerade veckor
från vecka 30 till och med söndagen före $r$. Låt $\mathbf X_r$ ha raderna
$\mathbf{x}_t^{\mathsf T}$ och låt $\mathbf y_r$ innehålla motsvarande
observationer. Koefficienterna bestäms genom
\begin{equation}
  \widehat{\boldsymbol\beta}_r=
  \underset{\boldsymbol\beta\geq0}{\operatorname{argmin}}
  \left\{
  \frac{\lVert\mathbf y_r-\mathbf X_r\boldsymbol\beta\rVert_2^2}{c^2}
  +\lambda\lVert\boldsymbol\beta-\boldsymbol\beta_0\rVert_2^2
  \right\}.
  \label{eq:ridge}
\end{equation}
Här är $c$ roten ur medelkvadraten av de historiska veckovärdena för platsen,
beräknad över vecka 30--23. Skalningen gör att samma $\lambda$ har jämförbar
innebörd för platser med olika fallnivåer. Icke-negativitetsvillkoret hindrar
modellen från att skapa svårtolkade kompensationer mellan positiva och
negativa säsongsvikter. Koefficienterna tvingas däremot inte summera till ett;
modellen kan därför anpassa den totala nivån uppåt eller nedåt.

För referensdatum $r$ beräknas de fyra prognoserna som
\begin{equation}
  \widehat y_{r+h}=\mathbf{x}_{r+h}^{\mathsf T}
  \widehat{\boldsymbol\beta}_r,\qquad h=0,1,2,3.
  \label{eq:horizons}
\end{equation}
Samma koefficientvektor används för alla fyra horisonter men skattas om inför
varje ny prognosomgång. Om den aktuella ISO-veckan är 53, vilken saknas i de
flesta historiska säsonger, används medelvärdet av vecka 52 och vecka 1 som
historisk kovariat.

\section*{Regularisering och kalibrering}

Ekvation~\eqref{eq:ridge} kan tolkas som en maximum-a-posteriori-skattning i
en Gaussisk regressionsmodell med en normalfördelad prior centrerad kring de
lika vikterna i ekvation~\eqref{eq:prior}, trunkerad till den icke-negativa
ortanten. Stora värden på $\lambda$ håller modellen nära säsongsmedelvärdet,
medan små värden tillåter en snabbare anpassning till den aktuella säsongen.

Ett gemensamt regulariseringsvärde används för samtliga platser. Kandidaterna
var
\[
  0.01,\ 0.03,\ 0.1,\ 0.3,\ 1,\ 3,\ 10,\ 30,\ 100,\ 300,\ 1000.
\]
De jämfördes med rullande prognosursprung i tidigare fullständiga säsonger.
För varje kalibreringssäsong användes endast ännu tidigare säsonger som
kovariater, och minst tre kovariatsäsonger krävdes. Urvalskriteriet var
genomsnittligt absolutfel över platser, prognosdatum och de fyra horisonterna.
Det bästa gemensamma värdet var
\[
  \boxed{\lambda=1}.
\]
Säsongen 2025/26 hölls helt utanför detta val och användes endast för den
slutliga retrospektiva utvärderingen.

\section*{Retrospektiv kontroll och operationell användning}

På 2025/26 gav den gemensamma regulariseringen ett totalt MAE på 96,7 fall,
jämfört med 98,5 när $\lambda$ valdes separat per plats. MAE var 229,8 för
Sverige, 27,6 för Skåne och 32,8 för Västra Götaland. Resultatet stöder den
enklare gemensamma specifikationen men ska tolkas som kontroll på en enda
influensasäsong, inte som ett generellt prestandalöfte.

I drift läser modellen den dataversion som är tillgänglig för veckans
prognosomgång, tar bort observationer efter söndagen före referensdatumet,
identifierar fullständiga tidigare säsonger och skattar
$\widehat{\boldsymbol\beta}_r$ separat för varje plats. Utdata består av tolv
icke-negativa punktprognoser: tre platser gånger fyra horisonter. Modellen
kvantifierar inte prognososäkerhet; probabilistiska prognoser skapas i hubben
av den separata ensemblemodellen.

Modellens viktigaste begränsning är antagandet att framtida säsongsförlopp kan
representeras av historiska säsonger. Nya virusvarianter, ändrad provtagning
eller strukturella förändringar i rapporteringen kan därför ge systematiska
fel. De skattade säsongsvikterna, deras summa samt prognosfelen bör följas under
säsongen, men $\lambda$ ska inte ändras med data från den säsong som samtidigt
utvärderas.

\end{document}
